\section*{Hoofdstuk \ref{ch:b1:precipitation} --- Neerslagvorming en oplosbaarheid}

\begin{solution}{exo:b1:precipitation:1}
\ce{BaSO4(s) <=> Ba^2+ + SO4^2-}, $K_s = [\ce{Ba^2+}][\ce{SO4^2-}] =
10^{-9.97}$. \ce{CaF2(s) <=> Ca^2+ + 2F-}, $K_s = [\ce{Ca^2+}][\ce{F-}]^2 =
10^{-9.84}$. \ce{Ag2CrO4(s) <=> 2Ag+ + CrO4^2-}, $K_s =
[\ce{Ag+}]^2[\ce{CrO4^2-}] = 10^{-11.95}$. \ce{Fe(OH)3(s) <=> Fe^3+ + 3OH-},
$K_s = [\ce{Fe^3+}][\ce{OH-}]^3 = 10^{-38.55}$. Voor $K_s = \num{2.0e-13}$
is $\mathrm{p}K_s = 12.70$.
\end{solution}

\begin{solution}{exo:b1:precipitation:2}
$s = \sqrt{K_s} = 10^{-12.27/2} = \qty{7.3e-7}{mol/L}$, dat is
$\num{7.3e-7} \times 187.8 = \qty{1.4e-4}{g/L} = \qty{0.14}{mg/L}$.
\end{solution}

\begin{solution}{exo:b1:precipitation:3}
Na het mengen (volume verdubbeld) is $[\ce{Ba^2+}] = \qty{1.0e-3}{mol/L}$
en $[\ce{SO4^2-}] = \qty{1.0e-4}{mol/L}$: $Q_s = \num{1.0e-7} > K_s =
\num{1.1e-10}$, bariumsulfaat slaat neer. In het tweede geval is
$[\ce{SO4^2-}] = \num{1.0e-7}$ en $Q_s = \num{1.0e-10} < K_s$: geen
neerslag, al is het op het nippertje.
\end{solution}

\begin{solution}{exo:b1:precipitation:4}
In water is $s = 10^{-9.97/2} = \qty{1.0e-5}{mol/L}$. In sulfaat van
\qty{0.010}{mol/L} (\cref{prop:b1:precipitation:common-ion}) is
$s = 10^{-9.97}/0.010 =
\qty{1.1e-8}{mol/L}$, ongeveer duizend keer minder (factor 970).
\end{solution}

\begin{solution}{exo:b1:precipitation:5}
In water is $s = (10^{-9.84}/4)^{1/3} = \qty{3.3e-4}{mol/L}$, of
\qty{26}{mg/L}. In calciumchloride van \qty{0.10}{mol/L} is
$[\ce{Ca^2+}] \approx 0.10$ en $[\ce{F-}] = 2s$, dus
$s = \frac12\sqrt{10^{-9.84}/0.10} = \qty{1.9e-5}{mol/L}$. De benadering
$s \ll 0.10$ klopt met vier grootteordes marge.
\end{solution}

\begin{solution}{exo:b1:precipitation:6}
Begin: $\mathrm{pH} = 14.00 - \frac12(11.25 + \log 0.050) = 14.00 - 4.97 =
9.03$. Bij \qty{99}{\percent} is $[\ce{Mg^2+}] = \num{5.0e-4}$ en
$\mathrm{pH} = 14.00 - \frac12(11.25 - 3.30) = 10.03$. Op een pH-as
bestaat \ce{Mg(OH)2} boven 9.03; daaronder is het magnesium volledig
\ce{Mg^2+} op \qty{0.050}{mol/L}.
\end{solution}

\begin{solution}{exo:b1:precipitation:7}
Zilverbromide verschijnt eerst, bij $[\ce{Ag+}] = 10^{-12.27 + 2} =
10^{-10.27}$, zilverchloride bij $10^{-7.75}$. Wanneer zilverchloride
verschijnt, is $[\ce{Br-}] = 10^{-12.27}/10^{-7.75} = 10^{-4.52} =
\qty{3.0e-5}{mol/L}$, \qty{0.30}{\percent} van het bromide. Met het
gebruikelijke criterium van \qty{0.1}{\percent} is de scheiding niet
aanvaardbaar: de twee $\mathrm{p}K_s$-waarden verschillen maar 2.5.
\end{solution}

\begin{solution}{exo:b1:precipitation:8}
Zilverchromaat ontstaat wanneer $[\ce{Ag+}]^2 \times \num{5.0e-3} = 10^{-11.95}$,
$[\ce{Ag+}] = \qty{1.5e-5}{mol/L}$. Dan is $[\ce{Cl-}] = 10^{-9.75}/\num{1.5e-5}
= \qty{1.2e-5}{mol/L}$, een fractie \num{2.4e-4} (\qty{0.024}{\percent})
van het chloride: de rode kleur verschijnt wanneer het chloride praktisch
volledig neergeslagen is.
\end{solution}

\begin{solution}{exo:b1:precipitation:9}
$[\ce{CO3^2-}] = [\ce{HCO3-}]K_{a2}/[\ce{H3O+}] = [\ce{HCO3-}]\,10^{\mathrm{pH}
- \mathrm{p}K_{a2}}$, dus $K_s = [\ce{Ca^2+}][\ce{HCO3-}]\,10^{\mathrm{pH} -
\mathrm{p}K_{a2}}$, en dat is de gevraagde betrekking. Met $[\ce{Ca^2+}] =
[\ce{HCO3-}] = s$: $s^2 = 10^{-8.30 + 10.33 - 8.30} = 10^{-6.27}$,
$s = \qty{7.3e-4}{mol/L}$ (\qty{73}{mg/L} \ce{CaCO3}), tegen
$10^{-4.15} = \qty{7.1e-5}{mol/L}$ als het carbonaation geen base was.
\end{solution}

\begin{solution}{exo:b1:precipitation:10}
1.~Optellen met \ce{CO2(g) <=> CO2(aq)}: $K' = K K_H = 10^{-4.34 - 1.47} =
10^{-5.81}$.
2.~$[\ce{Ca^2+}] = s$ en $[\ce{HCO3-}] = 2s$, dus $4s^3 = K'p$. Bij
\qty{0.010}{bar}: $s = (10^{-5.81} \times 0.010/4)^{1/3} =
\qty{1.6e-3}{mol/L}$, \qty{157}{mg/L}. Onder buitenlucht:
$s = \qty{5.5e-4}{mol/L}$, \qty{55}{mg/L}. (Controle: onder bodemlucht is
$[\ce{CO2(aq)}] = 10^{-3.47}$ en $\mathrm{pH} = 6.37 + \log(3.14\times
10^{-3}/3.39 \times 10^{-4}) = 7.34$: waterstofcarbonaat overheerst
inderdaad.)
3.~Water dat onder de hoge koolstofdioxidedruk van de bodem met calciet
verzadigd is, bereikt het plafond van de grot, waar de druk lager is:
het verliest koolstofdioxide, $Q > K'$, en calciet slaat neer waar de
druppel hangt, ring na ring.
% ledger: air:CO2, dfg:CO2_g, dfg:CO2_ao (log KH computed in figdata/bachelor-1/aqueous-constants.py)
\end{solution}

\begin{solution}{exo:b1:precipitation:11}
1.~Begin bij $\mathrm{pH} = 14.00 - \frac12(16.38 - 2.00) = 6.81$.
Volledig heroplossen wanneer $[\ce{Zn(OH)4^2-}] = 10^{-2} = 10^{-1.93}[\ce{OH-}]^2$,
$[\ce{OH-}] = 10^{-0.035}$, $\mathrm{pH} = 13.97$. Het hydroxide bestaat
tussen pH 6.81 en 13.97.
2.~$\mathrm{pH} = 14.00 - 14.45/4 = 10.39$, $s_{\min} = 2 \times
10^{-(16.38 + 1.93)/2} = \qty{1.4e-9}{mol/L}$ (in dit model met twee
deeltjes).
3.~Een rechte met richtingscoëfficiënt $-2$ van $(6.81, -2)$ omlaag naar
het minimum bij ongeveer $(10.39, -8.9)$, daarna richtingscoëfficiënt
$+2$ omhoog tot $(13.97, -2)$; de echte kromme wordt afgerond door de
andere hydroxodeeltjes (figuur in \cref{sec:b1:precipitation:amphoteric}).
\end{solution}

\begin{solution}{exo:b1:precipitation:12}
1.~$\ce{Al(OH)3(s) + OH- <=> Al(OH)4-}$, $K = 10^{-1.23}$; al het
aluminium is opgelost wanneer $10^{-3} \leqslant K[\ce{OH-}]$, dus
$[\ce{OH-}] \geqslant 10^{-1.77}$, $\mathrm{pH} \geqslant 12.23$.
2.~$[\ce{Fe^3+}] = 10^{-38.55}/10^{-3 \times 1.77} = 10^{-33.2}$: het
ijzer blijft volledig als \ce{Fe(OH)3}; filtreren scheidt het ijzer
(vaste stof) van het aluminium (filtraat).
3.~Magnesiumhydroxide is niet amfoteer: in een overmaat hydroxide blijft
het vast, samen met ijzer(III)hydroxide, en er wordt niets gescheiden.
\end{solution}

\begin{solution}{pb:b1:precipitation:1}
\textbf{1.} \ce{Fe(OH)3(s) <=> Fe^3+ + 3OH-}, $K_s = [\ce{Fe^3+}][\ce{OH-}]^3$;
\ce{Zn(OH)2(s) <=> Zn^2+ + 2OH-}, $K_s = [\ce{Zn^2+}][\ce{OH-}]^2$;
\ce{Mg(OH)2(s) <=> Mg^2+ + 2OH-}, op dezelfde manier.
\textbf{2.} De constanten hebben niet dezelfde vorm (macht 3 of 2 van
$[\ce{OH-}]$) en de drie ionen hebben verschillende concentraties; alleen
de pH bij het begin van de neerslagvorming kan vergeleken worden.
\textbf{3.} Bij het begin is $c[\ce{OH-}]^n = K_s$, dus $n\log[\ce{OH-}] =
-\mathrm{p}K_s - \log c$ en $\mathrm{pH} = \mathrm{p}K_e +
\log[\ce{OH-}] = \mathrm{p}K_e - \frac1n(\mathrm{p}K_s + \log c)$.
\textbf{4.} $14.00 - \frac13(38.55 - 3.00) = 2.15$.
\textbf{5.} $14.00 - \frac12(16.38 - 2.30) = 6.96$.
\textbf{6.} $14.00 - \frac12(11.25 - 1.70) = 9.22$.
\textbf{7.} IJzer(III) (2.15), dan zink (6.96), dan magnesium (9.22). Bij
pH 2.0 is voor ijzer $Q_s = 10^{-3} \times 10^{-36} = 10^{-39} < 10^{-38.55}$,
en de andere zijn nog verder van neerslaan: er is niets vast.
\textbf{8.} \qty{1.0e-6}{mol/L}.
\textbf{9.} $14.00 - \frac13(38.55 - 6.00) = 3.15$.
\textbf{10.} $\log[\ce{Fe^3+}] = -38.55 + 3(14.00 - \mathrm{pH}) = 3.45 -
3\,\mathrm{pH}$: richtingscoëfficiënt $-3$.
\textbf{11.} Tot het begin van de neerslagvorming van zinkhydroxide, pH
6.96.
\textbf{12.} Van 3.15 tot 6.96.
\textbf{13.} \ce{Fe^3+ + OH- <=> FeOH^2+}, $\beta_1 =
[\ce{FeOH^2+}]/([\ce{Fe^3+}][\ce{OH-}]) = 10^{11.82}$.
\textbf{14.} Beide zijn gelijk wanneer $[\ce{OH-}] = 10^{-11.82}$, bij pH
2.18: \ce{Fe^3+} overheerst daaronder, \ce{FeOH^2+} daarboven.
\textbf{15.} $10^{11.82 + 3.15 - 14.00} = 10^{0.97} = 9.3$: bij pH 3.15 is
het opgeloste ijzer 10.3 keer het vrije \ce{Fe^3+}, \qty{1.0e-5}{mol/L}:
er wordt maar \qty{99.0}{\percent} verwijderd. De eerste schatting was
niet veilig.
\textbf{16.} $[\ce{Fe}]_{\mathrm{opgelost}} = 10^{3.45 - 3\,\mathrm{pH}}\,
(1 + 10^{\mathrm{pH} - 2.18})$.
\textbf{17.} $\log[\ce{FeOH^2+}] = 11.82 + \log[\ce{Fe^3+}] + \mathrm{pH} -
14.00 = 1.27 - 2\,\mathrm{pH}$: richtingscoëfficiënt $-2$.
\textbf{18.} $1.27 - 2\,\mathrm{pH} \approx -6.00$ geeft 3.64; met de
term van \ce{Fe^3+} erbij is $10^{3.45 - 10.92}(1 + 10^{1.46}) =
\num{1.0e-6}$ bij $\mathrm{pH} = 3.64$.
\textbf{19.} $[\ce{Zn^2+}] = \num{5.0e-5}$: $14.00 - \frac12(16.38 - 4.30) =
7.96$. Magnesiumhydroxide begint bij 9.22: nee.
\textbf{20.} \ce{Zn(OH)2(s) + 2OH- <=> Zn(OH)4^2-}, $K = \beta_4 K_s =
10^{14.45 - 16.38} = 10^{-1.93}$.
\textbf{21.} $[\ce{OH-}]^2 = \num{5.0e-3}/10^{-1.93}$, $\mathrm{pH} =
13.81$. De operator mag niet doorschieten met base: boven pH 9.2 slaat
het magnesium samen met het zink neer, en ver daarboven gaat het zink
terug in oplossing.
\textbf{22.} $\num{1.0e-3} \times 1000 \times 0.999 \times 106.8 =
\qty{107}{g}$ per kubieke meter.
\textbf{23.} $\num{5.0e-3} \times 1000 \times 0.99 \times 99.4 =
\qty{492}{g}$ per kubieke meter.
\textbf{24.} Wordt de pH in één keer boven 6.96 gebracht, dan slaan beide
hydroxiden neer in één slib: het zink gaat verloren in het ijzerafval,
en het afval is met zink verontreinigd. Stoppen tussen de drempels
verwijdert alleen het ijzer.
\textbf{25.} \textbf{Van pH 3.64 tot pH 6.96}: onder 3.64 is nog meer dan
\qty{0.1}{\percent} van het ijzer opgelost, grotendeels als \ce{FeOH^2+};
boven 6.96 ontstaat zinkhydroxide. Wie \ce{FeOH^2+} verwaarloost, legt
de ondergrens bij 3.15.
\end{solution}
